\section{Numerical experiments}
\label{sec:ct-experiments}

We evaluate SFRAMES on spectrally constrained low-rank reconstruction from
real low-dose computed tomography (CT) measurements. The experiments examine
how sampling angular views affects progress at a fixed projection budget,
and compare SFRAMES with SPRING--SARAH, ProxSARAH and PALM. Increasing the
number of training views enlarges the finite sum while preserving the factor
dimensions, allowing us to compare minibatch and full-batch updates under
both measurement densities.
Evaluation uses \(14\) 2DeteCT slices from \(6\) specimen/material groups, with \(N\in\{360,1{,}080\}\) training views and \(360\) held-out views per slice.

\subsection{Spectrally constrained CT reconstruction}
\label{subsec:ct-model}

An attenuation image describes how strongly material reduces transmitted
X-ray intensity per unit distance. At each angular position, or view, the
source illuminates a line of detector bins along a fan of rays. Under the
Beer--Lambert model, the negative logarithm of normalized transmitted
intensity approximates the line integral of attenuation along a ray
\cite{detect}. We represent attenuation by an image
\(X\in\mathbb R^{n\times n}\), with \(n=512\). Let
\(\ell_{ij,pq}\) be the length in millimeters of ray \(j\) of view \(i\)
inside pixel \((p,q)\). The forward map for that view is
\begin{equation}
 P_i\colon\mathbb R^{n\times n}\longrightarrow\mathbb R^d,
 \qquad [P_i(X)]_j=\sum_{p,q}\ell_{ij,pq}X_{pq},\qquad d=956.
 \label{eq:ct-ray}
\end{equation}
The scanner geometry fixes the weights. The adjoint \(P_i^\ast\), called
backprojection, distributes detector residuals over the crossed pixels
using the same weights.

We fit factors \(U,V\in\mathbb R^{n\times r}\), with \(r=64\), so that
\(X=UV^\top\) has rank at most \(r\). For processed measurements
\(y_i\in\mathbb R^d\) and a training-view set \(\mathcal I\) of size
\(N\), define
\begin{equation}
 f(U,V)=\frac1N\sum_{i\in\mathcal I}q_i(U,V),\qquad
 q_i(U,V)=\frac{\|P_i(UV^\top)-y_i\|_2^2}{2d}.
 \label{eq:ct-loss}
\end{equation}
We bound each factor's largest singular value and impose entrywise
nonnegativity through
\begin{equation}
 \begin{split}
 C&=\{(U,V):\|U\|_{\rm op}\le R_U,\ \|V\|_{\rm op}\le R_V\},\\
 D&=\{(U,V):U\ge0,\ V\ge0\}.
 \end{split}
 \label{eq:ct-sets}
\end{equation}
With \(x=(U,V)\), \(T=\operatorname{Id}\) and \(g=\iota_D\), the
problem is
\begin{equation}
 \min_{x\in C,\,Tx\in D} f(x)
 \;=\;\min_{x\in C}\{f(x)+g(Tx)\}.
 \label{eq:ct-problem}
\end{equation}
Here \(\iota_D\) is zero on \(D\) and \(+\infty\) elsewhere.
Nonnegative factors ensure a nonnegative reconstructed attenuation image.
The set \(C\) is compact and convex, and strictly positive sufficiently
small factors lie in the interiors of both \(C\) and \(D\). Each component
has a Lipschitz gradient on a neighborhood of \(C\). A smooth compactly
supported extension equal to the component on that neighborhood gives a
globally Lipschitz gradient while preserving the values and gradients at
every queried factor pair.

For a uniform subset \(B\subset\mathcal I\), a joint stochastic gradient is
\begin{equation}
 E_B=\frac{1}{|B|d}\sum_{i\in B}P_i^\ast\bigl(P_i(UV^\top)-y_i\bigr),
 \qquad \widehat\nabla f_B(U,V)=(E_BV,E_B^\top U).
 \label{eq:ct-gradient}
\end{equation}
It is conditionally unbiased when \(B\) is drawn independently of the
current factors. Each selected view contributes all \(d\) detector bins.
One invocation uses a forward projection and an adjoint over the selected
views. The full-batch gradient uses \(B=\mathcal I\).

The orthant projection is \(\Pi_D(x)=x_+\), applied entrywise to both
factors. Thus the smoothed term is
\(g_{\beta}(x)=\operatorname{dist}_D(x)^2/(2\beta)\), with gradient
\((x-\Pi_D(x))/\beta\). The implemented SFRAMES direction and update are
\begin{equation}
 s_k=\operatorname{LMO}_C\!\left(m_k+
        \frac{x_k-\Pi_D(x_k)}{\beta_k}\right),\qquad
 x_{k+1}=(1-\gamma_k)x_k+\gamma_k s_k,
 \label{eq:ct-update}
\end{equation}
where \(m_0=\widehat\nabla f_{B_0}(x_0)\) and
\(m_k=(1-\alpha_k)m_{k-1}+\alpha_k\widehat\nabla f_{B_k}(x_k)\) for
\(k\ge1\). For a block direction \(G=A\operatorname{diag}(\sigma)B^\top\),
the spectral-ball LMO returns \(-RA_+B_+^\top\), with singular vectors
restricted to positive singular values. The two factors are handled
independently, using one thin SVD per factor.

\paragraph{Initialization and SFRAMES parameters.}
All methods start from common nonnegative rank-\(64\) factors \(U_0,V_0\),
obtained by an NNDSVD-style factorization of a sensitivity-normalized
backprojection of \(120\) training views, followed by a small positive
perturbation and a scale fit. Set \(R_U=2\|U_0\|_{\rm op}\) and
\(R_V=2\|V_0\|_{\rm op}\). Ten power iterations on the stacked calibration
operator \(P_J\) estimate \(\widehat\lambda\approx\|P_J\|_{\rm op}^2/(120d)\).
We use the initial block-curvature estimates to scale the smoothing prefactor:
\[
 L_U=\widehat\lambda\|V_0\|_{\rm op}^2,\qquad
 L_V=\widehat\lambda\|U_0\|_{\rm op}^2,\qquad
 \beta_0=c_\beta/\max(L_U,L_V).
\]
Both SFRAMES variants retain \((c_\beta,\gamma_0)=(16,0.1)\) from the original joint development search.
Minibatch SFRAMES draws \(64\) training views uniformly without replacement
at each update, independently across updates; full-batch SFRAMES uses all
\(N\) views. Both use
\[
 \gamma_k=\gamma_0(k+1)^{-3/4},\qquad
 \alpha_k=(k+1)^{-1/2},\qquad
 \beta_k=\beta_0(k+1)^{-1/4}.
\]

\paragraph{Evaluation criteria.}
Raw training loss is \(f(U,V)\). To evaluate a feasible reconstruction, we
repair each factor by
\[
 \mathcal R_R(A)=\frac{A_+}{\max\{1,\|A_+\|_{\rm op}/R\}},\qquad
 \widetilde X=\mathcal R_{R_U}(U)\mathcal R_{R_V}(V)^\top.
\]
Repaired training loss uses \(\widetilde X\) in~\eqref{eq:ct-loss}.
Held-out loss uses the same normalization over the \(360\) reserved views.
Relative image error is \(\|\widetilde X-X_{\rm ref}\|_F/\|X_{\rm ref}\|_F\),
where \(X_{\rm ref}\) is the aligned high-dose reconstruction. Feasibility
is measured by the orthant distance
\[
 d_D=\bigl(\|\min(U,0)\|_F^2+\|\min(V,0)\|_F^2\bigr)^{1/2}.
\]
The repair is used only for diagnostics.

\subsection{Results}
\label{sec:ct-results}

The comparison contains \(308\) complete evaluation runs on
\(14\) slices from \(6\) specimen/material groups,
with a common cap of \(69{,}120\) forward and \(69{,}120\) adjoint views,
four times the original budget. Table~\ref{tab:ct-results} reports quality and feasibility.
Median paired changes in repaired training loss relative to full-batch SFRAMES,
SPRING--SARAH, ProxSARAH and PALM, in that order, are
\((-7.9,-44.3,-19.4,-33.4)\%\) at \(N=360\) and \((-20.8,-40.2,-18.9,-53.7)\%\) at \(N=1{,}080\). Negative changes favor minibatch SFRAMES.
Paired changes are computed before aggregation and can differ from ratios of the table entries.
Across comparator/view-count combinations, minibatch SFRAMES is favored in
all \(6\) groups;
individual group changes range from \(-63.6\%\) to \(-5.8\%\).

\begin{table}[!htbp]
\centering
\caption{Reconstruction quality at the last recorded iterate within the extended budget. Entries take medians over seeds within slices, slices within specimen/material groups, then across groups. Raw and repaired are training losses; held-out loss and relative reference error use repaired factors. Lower is better. Bold gives the best value within each training-view count, including ties.}
\label{tab:ct-results}
\small
\setlength{\tabcolsep}{4pt}
\begin{tabular}{@{}llrrrrr@{}}
\toprule
$N$ & Method & Raw & Repaired & Held-out & Ref. error & $d_D$ \\
\midrule
360 & Minibatch SFRAMES & $\boldsymbol{0.0056562}$ & $\boldsymbol{0.0061095}$ & $\boldsymbol{0.0064943}$ & $0.3623$ & $0.05824$ \\
360 & Full-batch SFRAMES & $0.0063422$ & $0.0069884$ & $0.0071353$ & $\boldsymbol{0.3328}$ & $0.08823$ \\
360 & SPRING--SARAH & $0.04776$ & $0.04776$ & $0.047953$ & $0.6058$ & $\boldsymbol{0}$ \\
360 & ProxSARAH & $0.008613$ & $0.008613$ & $0.0086295$ & $0.4019$ & $\boldsymbol{0}$ \\
360 & PALM & $0.0099313$ & $0.0099313$ & $0.0099042$ & $0.4415$ & $\boldsymbol{0}$ \\
1{,}080 & Minibatch SFRAMES & $\boldsymbol{0.0058211}$ & $\boldsymbol{0.0062203}$ & $\boldsymbol{0.0063242}$ & $\boldsymbol{0.2830}$ & $0.05766$ \\
1{,}080 & Full-batch SFRAMES & $0.0078513$ & $0.0088956$ & $0.0088748$ & $0.3649$ & $0.1306$ \\
1{,}080 & SPRING--SARAH & $0.016491$ & $0.016491$ & $0.016495$ & $0.4834$ & $\boldsymbol{0}$ \\
1{,}080 & ProxSARAH & $0.0086672$ & $0.0086672$ & $0.0086554$ & $0.3982$ & $\boldsymbol{0}$ \\
1{,}080 & PALM & $0.014756$ & $0.014756$ & $0.014794$ & $0.5335$ & $\boldsymbol{0}$ \\
\bottomrule
\end{tabular}
\end{table}

Figures~\ref{fig:ct-trajectories-360} and~\ref{fig:ct-trajectories-1080}
show every recorded update against iterations, oracle calls and processed views.
The grey work-axis marker identifies the original \(17{,}280\)-view cap.
At the cap, every raw SFRAMES orthant distance exceeds \(1\%\) of the initial joint Frobenius norm. The largest absolute raw--repaired SFRAMES endpoint-loss gap is
\(20.19\%\) of repaired loss. The projection baselines remain in \(D\) throughout.

For the four joint targets, ordered as \((0.4,0.01),(0.4,0.001),(0.3,0.01),(0.3,0.001)\), the specimen-balanced attainment percentages, averaged equally over the two view counts, are \((3,0,3,0)\) for Minibatch SFRAMES; \((0,0,0,0)\) for full-batch SFRAMES; \((88,88,72,72)\) for SPRING--SARAH; \((100,100,94,94)\) for ProxSARAH; \((94,94,67,67)\) for PALM. Unreached targets remain right censored. Conditional on both methods attaining a target, specimen-balanced median SFRAMES/comparator work ratios range from \(1.06\) to \(6.68\) across \(5/32\) comparator/target/view-count combinations, each using \(2\text{--}3\) of \(42\) paired seed cases from \(1\) of \(6\) groups.

These are descriptive comparisons across the prespecified groups; slices and
oracle seeds within a group are repeated measurements. All retuned numerical baseline parameters are interior to their final grids.

\begin{figure}[p]
 \centering
 \includegraphics[width=\linewidth]{trajectories_360.pdf}
 \caption{Trajectories with \(360\) training views. Thin lines show
 individual runs, recorded at every update. Bold lines take median work
 and metric at the same iteration, first over seeds within slices and
 then over slices within groups and across groups. Bands show the middle \(50\%\) of group medians
 for SFRAMES. Training losses are normalized by their initial values. All loss axes are logarithmic.}
 \label{fig:ct-trajectories-360}
\end{figure}

\begin{figure}[p]
 \centering
 \includegraphics[width=\linewidth]{trajectories_1080.pdf}
 \caption{Trajectories with \(1{,}080\) training views, using the same
 metrics and summaries as Figure~\ref{fig:ct-trajectories-360}. The three
 projection baselines coincide at zero orthant distance.}
 \label{fig:ct-trajectories-1080}
\end{figure}

\clearpage
\subsection{Experimental protocol}
\label{subsec:ct-protocol}

We use Mode~1 low-dose fan-beam measurements from 2DeteCT \cite{detect};
the published Mode~2 high-dose reconstruction is the reference. Development
retains slices \texttt{04300} and \texttt{03150}. The prespecified evaluation
groups are Mix 1 (\texttt{00601}, \texttt{01201}); Mix 2 (\texttt{03041}, \texttt{03381}, \texttt{03611}); Mix 3 (\texttt{03791}, \texttt{04131}, \texttt{04931}); Almond (\texttt{05581}, \texttt{05611}); Lava stone (\texttt{05831}, \texttt{05861}); Walnut (\texttt{05731}, \texttt{05761}). Slices within each group are sections of the
same specimen/material acquisition, rather than independent specimens.

We discard the repeated final angle, sum adjacent detector cells, subtract
the similarly binned dark field, and divide by the mean flat field after
dark correction. Detector~1 acquisitions receive the published one-index
linear transmission shift with extrapolation before taking the negative
logarithm with floor \(10^{-6}\); Detector~2 uses zero shift. Reversing the
detector axis gives \(3{,}600\times956\) sinograms in the DeepInverse convention
\cite{deepinv}. The floor affects \(12\) in \texttt{05861}; it is inactive in the other evaluation acquisitions.
ASTRA~2.5.0 applies the FP32 CPU \texttt{line\_fanflat} projector \cite{astra},
while factor arithmetic is FP64. All acquisitions use source-to-origin and
origin-to-detector distances \(431.019989\) and \(97.980499\)~mm, detector
spacing \(0.1496\)~mm, and a square field of width \(116.528143\)~mm centered
at \((-0.113797,-0.113797)\)~mm. View \(i\) has angle \(\pi-2\pi i/3600\).
We vertically flip the \(1024\times1024\) reference, divide by its native
voxel width \(0.113797\)~mm, and average \(2\times2\) blocks.


For \(i\in\{0,\ldots,3599\}\), training selects \(i\bmod10=0\) for
\(N=360\) and \(i\bmod10\in\{0,3,6\}\) for \(N=1{,}080\).
Both reserve the \(360\) views with \(i\bmod10=9\). Calibration uses
the \(120\) training views with \(i\bmod30=0\), shared across both sizes.
Evaluation uses paired oracle seed labels \(202,203,204\), offset by
\(100{,}000\) for generators. PALM and full-batch SFRAMES run once per
slice/size and serve as shared deterministic references. Baseline development
uses labels \(110, 111, 112\), with PALM once per case. Calibration uses seed
\(110\) in development and \(201\) in evaluation.

All methods start from common calibrated factors in \(C\cap D\). We divide
\((P_J^\ast y_J)_+\) entrywise by \(P_J^\ast P_J\mathbf1\), with denominator
floor \(10^{-3}\) times its maximum. A nonnegative SVD-based factorization
gives rank-\(64\) factors, perturbed by \(10^{-3}\) times each factor's RMS
times uniform \([0,1)\) entries. A least-squares amplitude fit on calibration
views, bounded below by \(10^{-6}\), rescales both factors by its square root.
The numerical LMO retains singular values above
\(\epsilon_{\rm mach}\max(\operatorname{shape}G)\sigma_{\max}(G)\).

SPRING--SARAH \cite{spring} alternates factor updates with steps
\(a/L_U,a/L_V\). Each block independently refreshes with probability \(1/q\),
starting from a full gradient; otherwise it uses a paired \(64\)-view SARAH
difference. ProxSARAH \cite{proxsarah} uses joint steps \(a/(L_U+L_V)\),
snapshots of \(\min(N,256)\) views for \(s=\mathrm{mini}\) or all \(N\)
views for \(s=\mathrm{full}\), \(m\) recursive \(64\)-view updates, and relaxation
\(0.8\). Both terms of each recursive difference share sampled views.
PALM \cite{palm} uses full gradients and alternating steps \(a/L_U,a/L_V\).
Proximal steps project onto each spectral ball intersected with the orthant.
A feasible single-set projection is accepted directly; otherwise Dykstra uses
a primal--dual distance certificate, tolerance \(10^{-8}(k+1)^{-1.1}\),
at most \(2{,}000\) iterations, and checks at the first, every fifth and
final iteration, with floating-point allowance.
Feasibility rescaling tests the norm before division, including the zero-matrix case.

\begin{table}[!htbp]
\centering
\caption{Hyperparameter searches. Every row denotes the full Cartesian product of its listed values. SFRAMES settings are retained from the original joint search; the three baselines are retuned at the longer development horizon.}
\label{tab:ct-settings}
\small
\setlength{\tabcolsep}{4pt}
\begin{tabular}{@{}lll@{}}
\toprule
Method and parameter & Values tested & Selected \\
\midrule
\shortstack[l]{Minibatch SFRAMES\\$(c_\beta,\gamma_0)$} & \shortstack[l]{$c_\beta\in\{0.25,1,4,16,64\}$ \\ $\gamma_0\in\{0.05,0.1,0.16,0.25,0.5,0.8\}$} & $(16,0.1)$ \\
\shortstack[l]{Full-batch SFRAMES\\$(c_\beta,\gamma_0)$} & \shortstack[l]{$c_\beta\in\{0.25,1,4,16,64\}$ \\ $\gamma_0\in\{0.05,0.1,0.16,0.25,0.5,0.8\}$} & $(16,0.1)$ \\
SPRING--SARAH, $(a,q)$ & \shortstack[l]{$a\in\{0.7,1.1,1.6,2.3,3.4,5\}$ \\ $q\in\{8,16,32\}$} & $(2.3,16)$ \\
ProxSARAH, $(a,m,s)$ & \shortstack[l]{$a\in\{0.9,1.5,2.3,3.4,5\}$ \\ $m\in\{6,12,24,48,96,192\}$ \\ $s\in\{\mathrm{mini},\mathrm{full}\}$} & $(2.3,96,\mathrm{mini})$ \\
PALM, $a$ & \shortstack[l]{$a\in\{0.6,1.1,1.6,2.3,3.4,5\}$} & $2.3$ \\
\bottomrule
\end{tabular}
\end{table}

SFRAMES retains its joint \((c_\beta,\gamma_0)\) selection from
\(8{,}640\) views per direction and oracle label \(110\).
Figure~\ref{fig:ct-joint-selection} records that original search.
The baselines are retuned at \(34{,}560\) views per direction,
using complete Cartesian grids and extending selected edges to predefined
guards. Selection minimizes repaired-to-initial training-loss ratio, taking
seed medians within each development slice/size, then a median over the
four cases. Nonfinite development candidates are ineligible; operational
or projection failures are investigated. Evaluation groups, held-out views
and reference images do not enter parameter selection.

Evaluation permits \(69{,}120\)
forward and adjoint views each. Counters include every gradient evaluation,
refresh, snapshot and both terms of recursive differences. A gradient call
returns a joint or block gradient; each iteration updates both factors.
Processed views are \(W=\sum_{\ell=1}^G|B_\ell|\), where one view uses
one forward and one adjoint over all \(956\) bins. Diagnostic projections
are evaluated separately. Every update is recorded, and tables use the last
recorded iterate within budget. Existing trajectories continue from exact
optimizer/RNG checkpoints when parameters agree; new settings start from
the shared calibrated factors. Saved prefixes are preserved exactly.

Writing \(\mathcal R(x)\) for the pair of repaired factors, we fix joint
targets before expanded evaluation:
\[
 \frac{f(\mathcal R(x_k))}{f(x_0)}\le\rho,
 \qquad \frac{d_D(x_k)}{\sqrt{\|U_0\|_F^2+\|V_0\|_F^2}}\le\varepsilon,
 \quad \rho\in\{0.4,0.3\},\quad
 \varepsilon\in\{0.01,0.001\}.
\]
Both conditions must hold at the same recorded iterate. First-hit work uses
native observations, without interpolated crossings. Unreached targets are
right censored; work ratios use common attained pairs with coverage stated.
Attainment is averaged over seeds within slices and slices within groups,
then equally over groups and the two view counts. Quality and paired-effect
summaries instead take medians at each level, with groups equally weighted.

\begin{figure}[!htbp]
 \centering
 \includegraphics[width=0.9\linewidth]{joint_selection.pdf}
 \caption{Original SFRAMES joint development search at \(8{,}640\) views
 per direction. Cells give median repaired training loss relative to
 initialization over two development slices and both training-view counts,
 using oracle label \(110\). Here \(\beta_0=c_\beta/\max(L_U,L_V)\).
 Boxes mark the retained pairs; lower is better.}
 \label{fig:ct-joint-selection}
\end{figure}
